\documentclass[10pt,a4paper]{amsart}
\usepackage{amsmath,amssymb,amsthm,graphicx}
% Retained minimal preamble for the Stacks Project chapter sites-cohomology.tex (Cohomology on Sites).
% Verbatim excerpt of the chapter preamble: Xy-pic setup (native preamble lines 14-18) and the native
% theorem declarations with the mathematical macros (native preamble lines 154-217). Copied mechanically;
% nothing added, removed or redefined. The wrapper class and the amsmath/amssymb/amsthm/graphicx base
% packages are supplied by the assembler, so the class/externaldocument part of the native preamble is omitted.
% For commutative diagrams we use Xy-pic
\usepackage[all]{xy}

% We use 2cell for 2-commutative diagrams.
\xyoption{2cell}
% Theorem environments.
%
\theoremstyle{plain}
\newtheorem{theorem}[subsection]{Theorem}
\newtheorem{proposition}[subsection]{Proposition}
\newtheorem{lemma}[subsection]{Lemma}

\theoremstyle{definition}
\newtheorem{definition}[subsection]{Definition}
\newtheorem{example}[subsection]{Example}
\newtheorem{exercise}[subsection]{Exercise}
\newtheorem{situation}[subsection]{Situation}

\theoremstyle{remark}
\newtheorem{remark}[subsection]{Remark}
\newtheorem{remarks}[subsection]{Remarks}

\numberwithin{equation}{subsection}

% Macros
%
\def\lim{\mathop{\mathrm{lim}}\nolimits}
\def\colim{\mathop{\mathrm{colim}}\nolimits}
\def\Spec{\mathop{\mathrm{Spec}}}
\def\Hom{\mathop{\mathrm{Hom}}\nolimits}
\def\Ext{\mathop{\mathrm{Ext}}\nolimits}
\def\SheafHom{\mathop{\mathcal{H}\!\mathit{om}}\nolimits}
\def\SheafExt{\mathop{\mathcal{E}\!\mathit{xt}}\nolimits}
\def\Sch{\mathit{Sch}}
\def\Mor{\mathop{\mathrm{Mor}}\nolimits}
\def\Ob{\mathop{\mathrm{Ob}}\nolimits}
\def\Sh{\mathop{\mathit{Sh}}\nolimits}
\def\NL{\mathop{N\!L}\nolimits}
\def\CH{\mathop{\mathrm{CH}}\nolimits}
\def\proetale{{pro\text{-}\acute{e}tale}}
\def\etale{{\acute{e}tale}}
\def\QCoh{\mathit{QCoh}}
\def\Ker{\mathop{\mathrm{Ker}}}
\def\Im{\mathop{\mathrm{Im}}}
\def\Coker{\mathop{\mathrm{Coker}}}
\def\Coim{\mathop{\mathrm{Coim}}}

% Boxtimes
%
\DeclareMathSymbol{\boxtimes}{\mathbin}{AMSa}{"02}

%
% Macros for moduli stacks/spaces
%
\def\QCohstack{\mathcal{QC}\!\mathit{oh}}
\def\Cohstack{\mathcal{C}\!\mathit{oh}}
\def\Spacesstack{\mathcal{S}\!\mathit{paces}}
\def\Quotfunctor{\mathrm{Quot}}
\def\Hilbfunctor{\mathrm{Hilb}}
\def\Curvesstack{\mathcal{C}\!\mathit{urves}}
\def\Polarizedstack{\mathcal{P}\!\mathit{olarized}}
\def\Complexesstack{\mathcal{C}\!\mathit{omplexes}}
% \Pic is the operator that assigns to X its picard group, usage \Pic(X)
% \Picardstack_{X/B} denotes the Picard stack of X over B
% \Picardfunctor_{X/B} denotes the Picard functor of X over B
\def\Pic{\mathop{\mathrm{Pic}}\nolimits}
\def\Picardstack{\mathcal{P}\!\mathit{ic}}
\def\Picardfunctor{\mathrm{Pic}}
\def\Deformationcategory{\mathcal{D}\!\mathit{ef}}
\newenvironment{reference}{}{}
\newenvironment{slogan}{}{}
\newenvironment{history}{}{}
\usepackage{cleveref}
\crefname{originallemma}{original lemma}{original lemmas}
\Crefname{originallemma}{Original lemma}{Original lemmas}
\makeatletter
\newlabel{lemma-cech-map-into}{{03AS}{source}{}{sourcepointer-1}{}}
\newlabel{lemma-cech-map-into@cref}{{[originallemma][0][]03AS}{[1][0][]source}{}{}{}}
\makeatother
\crefname{originallemma}{original lemma}{original lemmas}
\Crefname{originallemma}{Original lemma}{Original lemmas}
\makeatletter
\newlabel{lemma-homology-complex}{{03AT}{source}{}{sourcepointer-2}{}}
\newlabel{lemma-homology-complex@cref}{{[originallemma][0][]03AT}{[1][0][]source}{}{}{}}
\makeatother
\crefname{originallemma}{original lemma}{original lemmas}
\Crefname{originallemma}{Original lemma}{Original lemmas}
\makeatletter
\newlabel{sites-modules-lemma-flat-resolution-of-flat}{{03EZ}{source}{}{sourcepointer-3}{}}
\newlabel{sites-modules-lemma-flat-resolution-of-flat@cref}{{[originallemma][0][]03EZ}{[1][0][]source}{}{}{}}
\makeatother
\title{The integral Cech complex tensored with any presheaf of rings stays exact in positive degrees}
\author{The Stacks Project Authors}
\date{stacks-project:sites-cohomology.tex:local-snapshot-20261011; GNU FDL 1.2 or later}
\begin{document}
\maketitle
\noindent\textit{Editorial note.} All delivered mathematical text is the author's own, copied from the pinned Stacks Project chapter file sites-cohomology.tex; no statement or proof step was rewritten, reordered or removed. The author offers a second proof route with the words 'Or you can use'; both routes are kept.\par\medskip
\noindent\textit{Editorial note.} The counted statement opens with the author's own slogan environment. The chapter's retained minimal preamble does not carry the project's no-op definitions of the display environments, so the documented no-op definitions of reference, slogan and history are added as a separate verbatim preamble span taken from the snapshot (sources/stacks-varieties/retained-preamble.tex lines 75-77, themselves the project's own wrapper definitions). Nothing was invented and no macro was redefined.\par\medskip
\noindent\textit{Editorial note.} Same-chapter references are located at their \textbackslash{}label lines inside the pinned chapter file; the cross-chapter pointer is located in the project's own tag list staged in the snapshot (sources/stacks-chow/tags.txt), at the verbatim line carrying its canonical tag.\par\medskip
\noindent\textit{Allowed standard prerequisites (editorial).}\begin{itemize}
\item The Cech chain complex of a family and the complex of modules it represents. Same chapter, Lemma labelled lemma-cech-map-into (source line 971, Stacks tag 03AS, counted earlier in this night's work); named in the counted statement at line 1167. Supplies the complex Z\_\{U, bullet\} being tensored; the exactness statement about it is the counted content.
\item Structure of the integral Cech complex at an object: a direct sum of complexes homotopic to Z in degree zero. Same chapter, Lemma labelled lemma-homology-complex (source line 1025, Stacks tag 03AT, already counted by the card DSH-alg-cech-homology-exact); invoked at proof lines 1174 and 1184. Supplies the description used in both branches of the proof; the tensor computation on top of it is retained in the counted span.
\end{itemize}\medskip
\ifdefined\hypertarget\hypertarget{sourcepointer-1}{}\fi
\noindent\textit{Source reference.} Original lemma 03AS: sources/stacks-sites-cohomology/pinned-original-source.tex, source lines 971--971. Same-chapter pointer named in the counted statement at line 1167 as the definition of the complex being tensored; the exactness argument is retained verbatim in the counted span.\par\medskip
\ifdefined\hypertarget\hypertarget{sourcepointer-2}{}\fi
\noindent\textit{Source reference.} Original lemma 03AT: sources/stacks-sites-cohomology/pinned-original-source.tex, source lines 1025--1025. Same-chapter pointer invoked at proof lines 1174 and 1184 for the shape of the integral Cech complex at an object; the tensor and exactness arguments are retained in the counted span.\par\medskip
\ifdefined\hypertarget\hypertarget{sourcepointer-3}{}\fi
\noindent\textit{Source reference.} Original lemma 03EZ: sources/stacks-chow/tags.txt, source lines 4474--4474. Cross-chapter standard prerequisite offered at proof lines 1181-1182 as the alternative route through flatness of the terms and of H\_0; the first route through homotopy is retained in the counted span.\par\medskip
\begin{lemma}
\label{lemma-complex-tensored-still-exact}
\begin{slogan}
The integral presheaf {\v C}ech complex is a flat resolution of the
constant presheaf of integers.
\end{slogan}
Let $\mathcal{C}$ be a category. Let
$\mathcal{U} = \{f_i : U_i \to U\}_{i \in I}$ be a family of morphisms
with fixed target such that all fibre products
$U_{i_0} \times_U \ldots \times_U U_{i_p}$ exist in $\mathcal{C}$.
Let $\mathcal{O}$ be a presheaf of rings on $\mathcal{C}$.
The chain complex
$$
\mathbf{Z}_{\mathcal{U}, \bullet}
\otimes_{p, \mathbf{Z}}
\mathcal{O}
$$
is exact in positive degrees. Here $\mathbf{Z}_{\mathcal{U}, \bullet}$
is the chain complex of Lemma \ref{lemma-cech-map-into}, and
the tensor product is over the constant presheaf of rings
with value $\mathbf{Z}$.
\end{lemma}

\begin{proof}
Let $V$ be an object of $\mathcal{C}$.
In the proof of Lemma \ref{lemma-homology-complex} we saw that
$\mathbf{Z}_{\mathcal{U}, \bullet}(V)$ is isomorphic as a complex
to a direct sum of complexes which are homotopic to $\mathbf{Z}$
placed in degree zero. Hence also
$\mathbf{Z}_{\mathcal{U}, \bullet}(V) \otimes_\mathbf{Z} \mathcal{O}(V)$
is isomorphic as a complex to a direct sum of complexes which are homotopic
to $\mathcal{O}(V)$ placed in degree zero.
Or you can use
Modules on Sites, Lemma \ref{sites-modules-lemma-flat-resolution-of-flat},
which applies since the presheaves $\mathbf{Z}_{\mathcal{U}, i}$ are flat,
and the proof of Lemma \ref{lemma-homology-complex} shows that
$H_0(\mathbf{Z}_{\mathcal{U}, \bullet})$ is a flat presheaf also.
\end{proof}

\end{document}
